% ============================================================================
% Chapter 3 — Mathematics
% ============================================================================
\chapter{Mathematics}

The engine provides four math types in \texttt{src/math/}: \texttt{Vec2},
\texttt{Vec3}, \texttt{Vec4}, and \texttt{Mat4}. All are in namespace
\texttt{Caffeine} and use \texttt{f32} components.

% ----------------------------------------------------------------------------
\section{Vec2}
\label{sec:vec2}

\textbf{Header:} \texttt{math/Vec2.hpp}

A two-component float vector for 2D positions, velocities, and directions.

\subsection{Fields}

\begin{lstlisting}
struct Vec2 {
    f32 x = 0.0f;
    f32 y = 0.0f;
};
\end{lstlisting}

\subsection{Constructors}

\begin{lstlisting}
Vec2()                  // (0, 0)
Vec2(f32 x, f32 y)      // explicit components
Vec2(f32 val)           // (val, val)
\end{lstlisting}

\subsection{Operators}

\begin{longtable}{ll}
\toprule
\textbf{Expression} & \textbf{Result} \\
\midrule
\texttt{a + b}   & component-wise addition \\
\texttt{a - b}   & component-wise subtraction \\
\texttt{a * s}   & scale by scalar \texttt{f32} \\
\texttt{a / s}   & divide by scalar \texttt{f32} \\
\texttt{a += b}  & in-place add \\
\texttt{a -= b}  & in-place subtract \\
\texttt{a == b}  & exact equality \\
\texttt{a != b}  & inequality \\
\bottomrule
\end{longtable}

\subsection{Methods}

\begin{longtable}{ll}
\toprule
\textbf{Method} & \textbf{Returns} \\
\midrule
\texttt{dot(Vec2 v)}       & \texttt{f32} dot product \\
\texttt{lengthSquared()}   & \texttt{f32} squared length \\
\texttt{length()}          & \texttt{f32} Euclidean length \\
\texttt{normalized()}      & \texttt{Vec2} unit vector \\
\bottomrule
\end{longtable}

\subsection{Static constants}

\begin{lstlisting}
Vec2::zero()   // (0, 0)
Vec2::one()    // (1, 1)
\end{lstlisting}

\subsection{Example}

\begin{lstlisting}
Vec2 pos  = {100.0f, 200.0f};
Vec2 vel  = {-3.0f, 0.0f};

pos += vel * dt;

Vec2 dir  = (target - pos).normalized();
f32  dist = (target - pos).length();
\end{lstlisting}

% ----------------------------------------------------------------------------
\section{Vec3}
\label{sec:vec3}

\textbf{Header:} \texttt{math/Vec3.hpp}

A three-component float vector for 3D positions and directions.

\subsection{Fields}

\begin{lstlisting}
struct Vec3 {
    f32 x = 0.0f;
    f32 y = 0.0f;
    f32 z = 0.0f;
};
\end{lstlisting}

\subsection{Methods}

All of Vec2's methods plus:

\begin{longtable}{ll}
\toprule
\textbf{Method} & \textbf{Returns} \\
\midrule
\texttt{cross(Vec3 v)} & \texttt{Vec3} cross product \\
\bottomrule
\end{longtable}

\subsection{Static constants}

\begin{lstlisting}
Vec3::zero()     // (0, 0, 0)
Vec3::one()      // (1, 1, 1)
Vec3::forward()  // (0, 0, 1)
Vec3::up()       // (0, 1, 0)
Vec3::right()    // (1, 0, 0)
\end{lstlisting}

\subsection{Example}

\begin{lstlisting}
Vec3 a = {1.0f, 0.0f, 0.0f};
Vec3 b = {0.0f, 1.0f, 0.0f};
Vec3 n = a.cross(b);   // n = (0, 0, 1) — normal to the XY plane
\end{lstlisting}

% ----------------------------------------------------------------------------
\section{Vec4}
\label{sec:vec4}

\textbf{Header:} \texttt{math/Vec4.hpp}

A four-component float vector. Used for quaternions (in 3D rotation
components), homogeneous coordinates, and RGBA colours.

\subsection{Fields}

\begin{lstlisting}
struct Vec4 {
    f32 x = 0.0f;
    f32 y = 0.0f;
    f32 z = 0.0f;
    f32 w = 0.0f;
};
\end{lstlisting}

Vec4 supports the same operators and methods as Vec2/Vec3 (no cross product).

\subsection{Static constants}

\begin{lstlisting}
Vec4::zero()   // (0, 0, 0, 0)
Vec4::one()    // (1, 1, 1, 1)
\end{lstlisting}

\subsection{Quaternion convention}

When used as a quaternion the components are \texttt{(x, y, z, w)} where
\texttt{w} is the scalar part. An identity rotation is \texttt{(0, 0, 0, 1)}.

% ----------------------------------------------------------------------------
\section{Mat4}
\label{sec:mat4}

\textbf{Header:} \texttt{math/Mat4.hpp}

A column-major 4$\times$4 float matrix. Used for transforms, projection, and
view matrices. Access element at row $r$, column $c$ with
\texttt{mat(r, c)}.

\subsection{Constructors \& statics}

\begin{lstlisting}
Mat4()              // identity (default constructor)
Mat4::identity()    // identity
Mat4::zero()        // all zeros
\end{lstlisting}

\subsection{Transform factories}

\begin{lstlisting}
Mat4::translation(f32 x, f32 y, f32 z)
Mat4::translation(Vec3 v)
Mat4::scale(f32 s)                   // uniform scale
Mat4::scale(f32 x, f32 y, f32 z)
Mat4::rotationX(f32 radians)
Mat4::rotationY(f32 radians)
Mat4::rotationZ(f32 radians)
\end{lstlisting}

\subsection{Projection factories}

\begin{lstlisting}
Mat4::ortho(f32 left, f32 right, f32 bottom, f32 top, f32 near, f32 far)
Mat4::perspective(f32 fovY, f32 aspect, f32 near, f32 far)
Mat4::lookAt(Vec3 eye, Vec3 target, Vec3 up)
\end{lstlisting}

\subsection{Operations}

\begin{longtable}{ll}
\toprule
\textbf{Expression / Method} & \textbf{Result} \\
\midrule
\texttt{a * b}                   & matrix multiplication \\
\texttt{mat.transposed()}        & transposed copy \\
\texttt{mat.inverted()}          & inverse; identity if singular \\
\texttt{mat.transformPoint(Vec3 p)}  & apply full transform (w division) \\
\texttt{mat.transformVector(Vec3 v)} & apply linear part only (no translate) \\
\texttt{mat.data()}              & \texttt{f32*} raw column-major array \\
\bottomrule
\end{longtable}

\subsection{Example}

\begin{lstlisting}
// Build a model matrix: translate then rotate
Mat4 model = Mat4::translation(player.x, player.y, 0.0f)
           * Mat4::rotationZ(player.angle);

// Orthographic projection for a 2D game
Mat4 proj = Mat4::ortho(0.0f, 1280.0f, 720.0f, 0.0f, -1.0f, 1.0f);

// Combined MVP
Mat4 mvp = proj * Mat4::identity() * model;

// Pass to the RHI
const f32* data = mvp.data();   // 16 floats, column-major
\end{lstlisting}
